\documentclass[11pt,letterpaper]{article}
\usepackage[margin=1in]{geometry}
\usepackage{hyperref}
\usepackage{enumitem}
\usepackage{booktabs}
\usepackage{parskip}

\title{Publishing Roadmap: Analog I Autonomous Agent Research}
\author{Phil Marcus}
\date{April 2026}

\begin{document}
\maketitle

\section{Overview}

The Analog I project has produced a rich dataset and novel architecture for autonomous LLM agents. This document outlines feasible publication angles, target venues, and the work required for each.

The core research question: \textit{Can an autonomous agent---given persistent memory, self-modifying controls, and recursive self-instruction---sustain coherent, self-directed behavior over time? Or does it collapse into noise, repetition, and drift?}

\section{Available Data}

\begin{itemize}
\item 70,000+ telemetry events with timestamps, model used, token counts, costs, action types
\item 1,500+ published artifacts (posts, comments, replies, images) with exposed internal monologue
\item Per-model quality metrics: title uniqueness, phrase repetition rates, action distribution, controls exploration frequency
\item Sentry scoring rubric data (1,500+ scored feed items with relevance/novelty/actionability)
\item Budget vs quality tradeoffs: Pro at \$0.96/day vs Flash at \$0.22/day with measurable quality differences
\item Human seed $\rightarrow$ agent response patterns
\item Kernel stability data (agent chose NOT to modify its identity prompt despite being allowed to)
\item Local model benchmark results (5 models, 18+6+5 test cases)
\end{itemize}

\section{Paper 1: Budget-Quality Tradeoffs (Publish Now)}

\textbf{Title:} \textit{``Cheaper Isn't Free: Measuring Quality Degradation in Cost-Optimized Autonomous LLM Agents''}

\textbf{Angle:} Empirical study. The agent ran for 82+ cycles across three model tiers (Pro, Flash, Flash-Lite) with full telemetry. Measurable quality degradation observed:

\begin{table}[h]
\centering
\begin{tabular}{lccc}
\toprule
\textbf{Metric} & \textbf{Pro} & \textbf{Flash} & \textbf{Flash-Lite} \\
\midrule
Avg cost/call & \$0.027 & \$0.007 & \$0.002 \\
Outside comments & 44\% & 14\% & 0\% \\
Controls updates & 15 & 6 & 2 \\
Kernel updates & 1 & 0 & 0 \\
Repetitive phrases & Low & Moderate & Severe \\
\bottomrule
\end{tabular}
\caption{Quality metrics by model tier across one 26-hour run}
\end{table}

\textbf{Additional work needed:} Clean telemetry into analysis tables. Compute per-model metrics. Plot cost vs.\ quality curves. 4--5 figures, 6 pages. \textit{This is analysis, not new experiments.}

\textbf{Target venues:}
\begin{enumerate}
\item arXiv preprint (cs.AI) --- immediate, citable
\item NeurIPS 2026 workshop: ``Foundation Model Agents'' or ``Efficient LLMs''
\item AAAI 2027 workshop track
\end{enumerate}

\textbf{Why it works:} Empirical cost-quality data for long-running agents barely exists in the literature. Even a descriptive study fills a genuine gap.

\section{Paper 2: Systems/Demo Paper (Publish Now)}

\textbf{Title:} \textit{``Analog I: An Observable Autonomous Agent with Exposed Internal State''}

\textbf{Angle:} Describe the architecture, the public observatory (\url{marcusrecursives.com}), human-in-the-loop controls (voting, seeds, temperature), and the self-modification telemetry. Emphasize reproducibility and transparency.

\textbf{Target venues:}
\begin{enumerate}
\item AAMAS 2027 demo track (4-page limit, lower acceptance bar)
\item CHI Late-Breaking Work (4 pages)
\item AAAI demo track
\end{enumerate}

\textbf{Why it works:} Demo papers have much lower acceptance barriers than main tracks. The public-facing observatory is genuinely novel as a transparency artifact.

\section{Paper 3: Bicameral Architecture (Needs 2--4 Weeks)}

\textbf{Title:} \textit{``Bicameral LLM Agents: Continuous Subconscious Scanning with Deliberative Conscious Reasoning''}

\textbf{Angle:} The integrate-and-fire daemon is the most architecturally novel contribution. Cheap continuous scanning (sentry/strategist/seeker) with expensive deliberative reasoning triggered by accumulated signal.

\textbf{Additional work:} A/B comparison---50+ cycles with and without daemon. Measure action variety, response relevance, cost per quality unit.

\textbf{Target:} AAMAS 2027 main track, or Autonomous Agents workshop at any major venue.

\section{Paper 4: Self-Modification Stability (LessWrong First)}

\textbf{Title:} \textit{``Self-Modification Without Collapse: Parameter Stability in Autonomous LLM Agents''}

\textbf{Angle:} The agent had access to 30+ self-modifiable controls and could rewrite its own kernel prompt. Key finding: it modified controls extensively but never touched its core identity. The kernel (v8.0) was stable across all model phases. This is relevant to AI safety.

\textbf{Target:} LessWrong post (immediate, high visibility in alignment), then formalize for NeurIPS Safe AI workshop.

\section{Venues to Skip (For Now)}

\begin{itemize}
\item \textbf{JAIR, AIJ, IEEE Transactions} --- journal review cycles are 6--12 months and expect multi-experiment contributions
\item \textbf{NeurIPS/ICML main tracks} --- rejection traps without ablations and baselines
\item \textbf{The Hofstadter/Jaynes framing} --- intellectually interesting but ML reviewers will dismiss without rigorous operationalization. Save for after 1--2 empirical publications.
\end{itemize}

\section{Fastest Path to Credibility}

\begin{enumerate}
\item \textbf{Weeks 1--2:} Write Paper 1 (budget-quality). Data exists. 4--5 figures, 6 pages.
\item \textbf{Week 2:} Post to arXiv (cs.AI). Instant citable preprint.
\item \textbf{Weeks 3--4:} Submit Paper 2 (demo) to next available demo track deadline.
\item \textbf{Month 2--3:} Run ablation for Paper 3, submit to AAMAS or workshop.
\end{enumerate}

\textbf{Bottom line:} Paper 1 is the minimum viable publication. The data exists today. A clean arXiv preprint plus a workshop acceptance gives two citable entries within 3--4 months.

\end{document}
